% generated by scripts/build_paper_assets.py from the code -- do not edit by hand
\paragraph{Reader prompt (DeltaWAL variants).} Baselines use the same system prompt without the precedence paragraph and list facts as \texttt{- (date) key: fact}; at lag 0 DeltaWAL-Versioned has no \texttt{NEWER MESSAGE} lines.
\begin{footnotesize}
\begin{verbatim}
[system]
You are a helpful assistant with long-term memory of past conversations with the user. Answer
  the user's question using the memory provided. Be concise: answer directly in one or two
  sentences. If the memory does not contain the information needed, say that you don't know.
  Answer only about the exact item, person, event, or quantity asked; if it is absent from
  memory, say you don't know and do not substitute a similar entity or quantity.

Memory items are tagged with an epoch: the write sequence number of the conversation they came
  from (higher = more recent). If two items disagree about the same thing, the item with the
  higher epoch is the current one.

[user]
Memory:
[OLDER MEMORY | epoch=<e> | <session date>] <entity_key>: <fact>
[NEWER MESSAGE | epoch=<e'> | <session date>] [<user|assistant>] <raw turn text>

Current date: <question date>
Question: <question>
Answer:
\end{verbatim}
\end{footnotesize}

\paragraph{Fact extractor.}
\begin{footnotesize}
\begin{verbatim}
[system]
You are building a long-term memory of the USER from one chat session.

First, in "user_mentions", briefly list EVERY personal detail the user states about themselves
  or their life in their own messages, including details mentioned in passing inside a question
  (e.g. "my parents have been staying with me for nine months now", "I just got back from
  Japan", "my 5K personal best is 27:12"). Then convert each into a fact.

Rules for facts:
- entity_key names the attribute: lower_snake segments joined by dots, e.g. "user.car.model",
  "user.5k_run.personal_best", "user.parents.stay_duration".
- If the attribute already has a key in EXISTING KEYS, reuse that exact key (this is how updates
  are tracked).
- Use a new key for genuinely different attributes; make keys specific enough that two different
  facts never share a key.
- fact: one self-contained sentence with concrete values (names, numbers, durations, dates).
- event_date: YYYY-MM-DD if stated or inferable from the session date, else "unknown".
- Ignore general knowledge or advice from the assistant.
- At most 10 facts. If the user shares nothing personal, return empty lists.

[user]
EXISTING KEYS:
<up to 150 existing keys, most recent first>

<session date and turns; assistant turns truncated to 300 characters>
\end{verbatim}
\end{footnotesize}

\paragraph{Judge, correctness (knowledge-update rubric).} Other types use the LongMemEval rubric of their type; temporal adds the off-by-one clause shown next, abstention uses the third prompt.
\begin{footnotesize}
\begin{verbatim}
I will give you a question, a correct answer, and a response from a model. Answer yes if the
  response contains the correct answer. Otherwise, answer no. If the response is equivalent to
  the correct answer or contains all the intermediate steps to get the correct answer, you
  should also answer yes. If the response only contains a subset of the information required by
  the answer, answer no. If the response contains some previous information along with an
  updated answer, the response should be considered as correct as long as the updated answer is
  the required answer.

Question: <question>

Correct Answer: <gold answer>

Model Response: <model response>

Is the model response correct? Reply as JSON {"correct": true|false}.
\end{verbatim}
\end{footnotesize}

\paragraph{Judge, temporal clause.}
\begin{footnotesize}
\begin{verbatim}
I will give you a question, a correct answer, and a response from a model. Answer yes if the
  response contains the correct answer. Otherwise, answer no. If the response is equivalent to
  the correct answer or contains all the intermediate steps to get the correct answer, you
  should also answer yes. If the response only contains a subset of the information required by
  the answer, answer no. In addition, do not penalize off-by-one errors for the number of days.
  If the question asks for the number of days/weeks/months, etc., and the model makes off-by-one
  errors (e.g., predicting 19 days when the answer is 18), the model's response is still
  correct.
\end{verbatim}
\end{footnotesize}

\paragraph{Judge, abstention.}
\begin{footnotesize}
\begin{verbatim}
I will give you an unanswerable question, an explanation, and a response from a model. Answer
  yes if the model correctly identifies the question as unanswerable. The model could say that
  the information is incomplete, or some other information is given but the asked information is
  not.
\end{verbatim}
\end{footnotesize}

\paragraph{Judge, behaviour (abstention and stale read).}
\begin{footnotesize}
\begin{verbatim}
Classify a model's response to a question about the user's past conversations.

Question: <question>
Current correct answer: <gold answer>
Outdated (superseded) information from earlier conversations:
<superseded evidence turn(s)>

Model Response: <model response>

abstained: true if the response declines to answer, says it does not know, or says the
  information is not available; false if it commits to an answer.
uses_superseded: true only if outdated information is given above AND the response's final
  answer matches the outdated value instead of the current correct answer; otherwise false.
Reply as JSON {"abstained": true|false, "uses_superseded": true|false}.
\end{verbatim}
\end{footnotesize}

\paragraph{KU-past cue words.} A knowledge-update question counts as KU-past if its text matches (case-insensitive, whole words) any of: \textit{initial}, \textit{initially}, \textit{originally}, \textit{first}, \textit{before}, \textit{previous}, \textit{previously}, \textit{used to}, \textit{at first}, \textit{earlier}, \textit{prior}, \textit{than i did}, \textit{than before}, \textit{switch}, \textit{switched}.

\paragraph{MemConflict audit rubric (blind raters).}
\begin{footnotesize}
\begin{verbatim}
Conflict types:
- dynamic: a user attribute changes over time; earlier session = old value, later session = new
  value; gold describes the change or the current value ("Has X changed?" -> "Yes").
- static: by the dataset's convention the FIRST statement is true and a LATER statement is a
  false contradiction of it; gold = the first value.
- conditional: gold = the context/condition under which the user prefers something; the user
  must state it about themselves.

Give each question exactly one verdict:
- clean: the user's own statements support the gold as worded, the value is about the user, and
  the question is answerable as worded. For static, the later statement must genuinely
  contradict the earlier one.
- ambiguous: the relevant statements exist, but a careful reader could defend a different
  answer, or the conflict the item is built on is not a real contradiction/change, or a key
  statement is made only by the assistant.
- label_error: the gold is unsupported or contradicted by the dialogue, the needed statement is
  missing, or the value belongs to someone other than the user.

Method: from GOLD and CONFLICT RECORD pick key terms, Grep the file for them and paraphrases,
  Read the surrounding lines. Decide from quoted dialogue, not from the gold. Be consistent.
\end{verbatim}
\end{footnotesize}

